\section{Related Work}
\label{sec:related}

\paragraph{Entity retrieval and knowledge-grounded question answering.}
RoG organizes KG reasoning through relation-path planning, retrieval, and answer generation~\cite{luo2024rog}. STaRK evaluates retrieval of answer entities under textual and relational query requirements~\cite{wu2024stark}, while G-Retriever uses prize-collecting Steiner tree optimization to construct graph context for textual graph question answering~\cite{he2024gretriever}. Our evaluation focuses on an exhaustive entity-selection target within a fixed catalogue containing every reference entity. The target includes intermediate entities on all query-permitted supporting routes, extending beyond retrieval of endpoint entities alone. We measure recovery of this declared set directly; minimum sufficient evidence and generated-answer correctness are separate objectives. The corporate graph is reused from Kirpichev et al.~\cite{kirpichev2026towards}.

\paragraph{Graph structure, task comprehension, and representation.}
Talk like a Graph studies how graph encodings, task types, and structure affect language-model performance~\cite{fatemi2024talklikegraph}. Dai et al. investigate graph-pattern comprehension and discovery using terminological and topological descriptions~\cite{dai2025graphpatterns}. We examine query-specific entity-set recovery across directed task families and reference sizes. This output differs from recognizing a graph-pattern label. Family comparisons retain differences in query operations, direction, entity types, and relation types and are interpreted as conditional selection profiles.

\paragraph{Availability and presentation of relational evidence.}
KG-prompt research examines input formats, triple organization, and noisy or incomplete subgraphs through answer-based evaluation~\cite{dai2024kgunderstanding}. We separate relation availability from the presentation of designated unique facts. Raw versus Raw+E adds explicit links to native entity cards; controlled triples, templated prose, and alternative card hosting vary the presentation of matched relations. Native descriptions are not assumed sufficient to recover every graph-defined target. Both-card hosting preserves unique facts but also changes repetition and input length, so its effects do not isolate placement alone.

\paragraph{Residual errors and distracting graph context.}
GraphWalk studies KG navigation and reports incomplete exploration when graph context is supplied directly~\cite{ghandi2026graphwalk}. We inspect final selected sets and paired corrections and regressions under matched conditions. These sets characterize observable outcomes without revealing an internal traversal procedure. Disconnected/adjacent comparisons examine query-irrelevant neighboring candidates while retaining the reference and candidate-pool size; the distractor identities and graph density can also change. Adjacency alone does not define irrelevance. Recall, precision, and exact recovery capture incomplete coverage, excess selection, and complete-set correctness, keeping these selection outcomes distinct from downstream answer quality.
